\documentclass[12pt]{article}
\usepackage[utf8]{inputenc}
\usepackage[T1]{fontenc}
\usepackage{hyperref}
\title{Kugu-Muminh Language Resources for K-12 Curriculum Development}
\author{Compiled by Claude Sonnet 5.5 (claude-sonnet-5-5) for Andr\'as Kornai}
\date{}
\begin{document}
\maketitle

\section{Introduction}
Kugu-Muminh (ISO 639-3 xmh; Glottocode kuku1283), also called Wik-Muminh or Kugu-Nganhcara, is a Paman (Pama-Nyungan) language of Cape York Peninsula, Queensland. Sources disagree on its status: Glottolog marks it dormant, while Endangered Languages cites a 1981 count of 31 speakers. Endangered Languages also reports no literacy tradition. Any curriculum would be a heritage and revitalisation one.

\section{Primary Educational Resources}
No primers, textbooks or school programmes were found in the search results.

\section{Linguistic and Reference Materials}
\begin{itemize}
\item Wikipedia, ``Kugu Nganhcara language'' (dialects and overview): \url{https://en.wikipedia.org/wiki/Kugu_Nganhcara_language}
\item Wikipedia, ``Wik languages'': \url{https://en.wikipedia.org/wiki/Wik_languages}
\item Glottolog entry: \url{https://glottolog.org/resource/languoid/id/kuku1283}
\item Endangered Languages Project entry: \url{https://endangeredlanguages.com/lang/3243}
\end{itemize}

\section{Religious and Community Resources}
None found in the search results.

\section{Digital and Audio Resources}
\begin{itemize}
\item Global Recordings Network page for xmh (Wik-Muminh): \url{https://globalrecordings.net/en/language/xmh}
\end{itemize}

\section{Curriculum Development Context}
The language is not known to be taught in schools. Documentation appears thin. Any revitalisation work should be led by the relevant Wik community and Aboriginal language centres. I did not verify the existence or contacts of such bodies in this search.
\end{document}